\chapter{Discussion and Conclusion}\label{ch:conclusion}

\section{Main findings}
The single-lead hierarchy achieved 98.07\% pooled six-class accuracy
on later accepted MITDB MLII beats from known patients. Equal-subject
accuracy is 98.03\%, with a subject-clustered 95\% interval of
96.29--99.25\%. The interval width was 2.95 percentage points, meeting the target of at most
three percentage points. The width was calculated from unrounded bounds. However, PAC recall is 70.63\%, and two subject clusters
remain below 90\% accuracy.

A limited check of complete INCART I01 and I02 records gives 94.11\%
binary accuracy and 92.85\% whole-system accuracy. Both records belong to
the same source patient. Binary precision falls to 64.22\% because 317
normal beats become false alarms. The lower precision shows that false alarms remain a major issue in this
external sample. The one-patient design limits the scope of this finding.

\section{Interpretation of the findings}
The high pooled accuracy shows that the system classified most supported
beats correctly in this temporal evaluation. It does not show uniform
performance across classes or patients. PAC recall was 70.63\%, compared
with more than 96\% for each of the other supported abnormal groups. The
113 PAC-to-PVC errors show that subtype confusion was a major source of
failure. Both the binary gate and subtype model therefore need improvement
for this group.

Conditional subtype accuracy measures how well the second model classifies
true supported abnormal beats when the gate is ignored. Whole-system accuracy
also includes the gate and the larger N group. The two percentages have
different denominators, so their difference is not a direct measure of the
cost of the gate. The class-level counts provide a clearer view of those
errors. Likewise, the wide conditional-subtype confidence interval shows
substantial variation across subjects despite high pooled accuracy.

The use of waveform and RR features follows earlier heartbeat-classification
research \cite{dechazal2004heartbeat}. Normal-reference reconstruction adds
another source of information, but this experiment does not isolate its
benefit. Differences in labels, leads, and patient splits also prevent a
direct accuracy ranking against the studies reviewed in
Chapter~\ref{ch:literature}.

The INCART check shows a further limitation. High sensitivity was accompanied
by 317 false alarms and binary precision of 64.22\%. Thus, detecting most
abnormal beats did not ensure high precision among the alerts. Changes in
recording conditions, lead configuration, and the proportion of abnormal
beats may contribute to this difference, but the one-patient sample cannot
separate their effects.

\section{Achievement of the objectives}
\begin{itemize}
\item The prototype integrates a frozen NSRDB autoencoder and two supervised
MLII ensembles, with the same local model-window preparation at training and use.
\item Pooled whole-system accuracy, the equal-subject lower confidence bound,
and the interval-width target are met in the corrected retrospective run.
\item The requirement that every supported abnormal group exceed 90\% recall
is not met because of PAC.
\item Generalisation across independent external patients, all supported
subtypes, and real AD8232 recordings is not established.
\end{itemize}

\section{Effect of the preprocessing revision}
The original supervised cache standardised complete MITDB recordings before
window extraction. Its scale statistics therefore included the test region,
and sensor chunks used a different scale. The corrected system filters and
standardises each four-second model window and requires both neighbouring
R peaks to be observed. Tests verify agreement between offline and streaming
feature construction when the same peak locations are supplied.

The original 96.86\% result is retained as a historical artifact, not the
headline result of the corrected system. The preprocessing changes, removal
of two edge beats per record, and refitting mean that the difference between
the two runs is not a controlled estimate of one component's causal benefit.

The autoencoder remains the original frozen checkpoint. Its decoder resizes
up-block tensors to their skip lengths and finally interpolates from 256 to
512 samples. Chapter~\ref{ch:architecture} reports these implemented dimensions. Altering this trained architecture would require a separate
training study.

\section{Limitations}
\begin{enumerate}
\item Earlier labelled ECG from the same patients is available during
supervised development. The result does not establish accuracy for patients whose data were absent
from model development.
\item The final segment was observed during earlier development. Correcting
the preprocessing does not restore test independence. Earlier development
choices may still make the results optimistic. Neither the bootstrap nor the
reported $p$-values corrects this bias.
\item The bootstrap estimates subject variation with the fitted model fixed.
It does not include retraining variability or extrapolation to a new cohort.
\item NSRDB is a nominally normal source cohort, with no beat-annotation
screening in this preparation. Its ECG1 anatomical lead is unspecified.
The two-epoch training record does not establish optimal convergence.
\item No controlled ablation quantifies the independent benefit of the
four autoencoder residuals.
\item Output labels are project groups. AFib includes flutter on normal-like
beat symbols; PVC includes ventricular escape beats; N includes atrial and
junctional escape beats. No AF episode detection is evaluated.
\item Unsupported abnormal beats enter binary testing but are excluded from
six-class scoring. The deployed gate has no validated unknown-class rejection.
\item Temporal and INCART scoring use annotated R peaks. Live peak detection
introduces errors beyond those metrics. The prepared window needs roughly
2.44 seconds of post-R signal and the next observed beat, plus acquisition
and computation delay.
\item INCART evaluation covers 5,416 beats but only one source patient and
N/PVC reference support. More beats from this same patient cannot supply
independent-patient evidence or validate absent classes.
\item The software requires the correct lead declaration, but cannot verify
physical electrode placement. Signal checks reject non-finite/flat inputs and
reported lead-off; they do not certify clinical signal quality.
\end{enumerate}

\section{Further research}
The next experiment should preregister a patient-disjoint development,
validation and test allocation; freeze preprocessing and thresholds before
external scoring; and recruit enough independent subjects for every target
class. PAC and external false alarms need particular attention. A controlled
comparison should evaluate the hierarchy with and without residual features,
and prospective sensor testing should measure R-peak accuracy, classification
errors, signal quality and end-to-end latency. Future work should treat the already viewed temporal and INCART records as
development data and reserve separate patients for final evaluation.

\section{Conclusion}
This study developed a reproducible single-lead ECG prototype and evaluated
its binary, subtype, and complete classification stages. The retrospective
results met the aggregate targets, but PAC recall and external false alarms
remain important weaknesses. It remains a research system, not a
validated clinical diagnostic device \cite{collins2024tripodai}.
